\section{底层模型对比：Opus 4.6 vs.\ GPT-5.3-Codex}

\subsection{Task-Completion Time Horizon}

衡量模型能力的一个重要指标是 \textbf{Task-Completion Time Horizon}（任务完成时间跨度）：给模型一个"人类专家需要 $T$ 小时完成"的任务，模型能以多大概率成功完成？

\begin{importantbox}{核心对比：任务完成能力差距显著}
\begin{itemize}
    \item \textbf{Opus 4.6}：在 50\% 成功率下，可处理长达 \textbf{12 小时}的任务
    \item \textbf{GPT-5.3-Codex}：在 50\% 成功率下，可处理长达 \textbf{5 小时 50 分钟}的任务
    \item 在 80\% 成功率下，两者差距明显缩小
\end{itemize}
这表明 Opus 4.6 在处理复杂、长周期任务时具有明显优势。
\end{importantbox}

\subsection{Token 经济性}

基准测试显示两者准确率接近，但 token 消耗差异巨大。Morph 的对比研究发现：

\begin{table}[H]
\centering
\caption{Token 消耗对比}
\begin{tabular}{lcc}
\toprule
\textbf{指标} & \textbf{Claude Code} & \textbf{Codex} \\
\midrule
相同任务 token 倍率 & 3.2--4.2$\times$ & 1$\times$（基准） \\
Figma 插件构建 & 6.2M tokens & 1.5M tokens \\
\bottomrule
\end{tabular}
\end{table}

\begin{warningbox}{Token 消耗的实际影响}
在相同订阅价格下，Claude Code 用户更可能触及 token 用量上限。如果你的工作流涉及大量代码生成，需要特别关注这一点。
\end{warningbox}

\subsection{本章小结}

Opus 4.6 在长任务上的可靠性显著高于 GPT-5.3-Codex，但 Claude Code 的 token 消耗远高于 Codex。选择时需要在"任务处理深度"和"token 效率"之间权衡。

